\subsection*{The complete nullvector equation without an $H^1$ assumption (NS-51)}

\paragraph{Classification and scope.}
The next statements are exact identities and proved implications, with a
named open arithmetic injectivity input. They concern the complete
fixed-window operator, not a finite head. The new regularity countermodel
rules out one generic inference, not injectivity of the arithmetic
operator. No new window or tail metric is computed. G2 and RH remain open.

Put $I_a=(-a,a)$, and let $E_a$ denote zero extension from $I_a$.
Write $V_a=\mathcal D(q_a)$ for the complete physical form domain:
\[
 V_a=\left\{f\in L^2(I_a):
  \int_{\mathbb R}\log(2+|\xi|)|\widehat{E_af}(\xi)|^2\,d\xi<\infty\right\}.
\]
Let $A_a$ be its self-adjoint form operator. These are the unshifted,
complete objects of Proposition~\ref{prop:v118-log-dirichlet}; no source
projection is imposed. Let $\mathfrak g$ be Suzuki's real even screw
function, normalized by $\mathfrak g(0)=0$ and, for $t>0$,
\begin{align}
 \mathfrak g(t)={}&-4(e^{t/2}+e^{-t/2}-2)
  +\sum_{n<e^t}\frac{\Lambda(n)}{\sqrt n}(t-\log n)
  -\frac t2\bigl(\psi(1/4)-\log\pi\bigr)\notag\\
 &-\frac14\left\{\Phi(1,2,1/4)
             -e^{-t/2}\Phi(e^{-2t},2,1/4)\right\}.
 \label{eq:ns51-screw}
\end{align}
Here $\Phi(z,s,v)=\sum_{j\ge0}z^j/(j+v)^s$ for $|z|<1$,
with its convergent value at $(1,2,1/4)$. Including a prime-power
endpoint instead makes no difference in this formula, since its ramp
vanishes at that endpoint. Suzuki's distributional identity is
$A_a\varphi=(-\mathfrak g''*E_a\varphi)|_{I_a}$ for
$\varphi\in C_c^\infty(I_a)$~\cite[Sections 2.5 and 3]{Suzuki2026}.
All convolutions below are local convolutions with a compactly supported
function or distribution; no global temperedness of $\mathfrak g$ is
assumed.

\begin{proposition}[Affine screw potential criterion on the complete domain]
\label{prop:ns51-affine-potential}
For $f\in L^2(I_a)$ set
\[
 F_f(x)=\int_{-a}^a\mathfrak g(x-y)f(y)\,dy,\qquad
 U_f(x)=\int_{-a}^a\mathfrak g'(x-y)f(y)\,dy.
\]
Then $F_f\in C^1(\mathbb R)$ and $F_f'=U_f\in C(\mathbb R)$.
For $f\in V_a$ the following conditions are equivalent:
\begin{enumerate}[label=(\roman*)]
 \item $f\in\mathcal D(A_a)$ and $A_af=0$;
 \item $F_f$ is affine on $I_a$;
 \item $U_f$ is constant on $I_a$.
\end{enumerate}
Let $\Pi_{\rm aff}$ be the ordinary $L^2(I_a)$ projection onto
$\operatorname{span}\{1,x\}$ and $P_0$ the projection onto mean-zero
functions. The operators
\[
 \mathcal C_a f=(I-\Pi_{\rm aff})F_f|_{I_a},\qquad
 \mathcal T_a f=P_0 U_f|_{I_a}
\]
are compact on $L^2(I_a)$, and
\begin{equation}\label{eq:ns51-kernel-reduction}
 \ker A_a=V_a\cap\ker\mathcal C_a=V_a\cap\ker\mathcal T_a.
\end{equation}
For even $f$ the affine potential is constant and $U_f=0$; for odd
$f$ it is $b x$ and $U_f=b$. The constant $b$ must not be discarded.
\end{proposition}
\begin{proof}
The screw function is locally absolutely continuous. Its derivative
has only a logarithmic singularity at zero and finitely many bounded
jumps on every compact interval. Indeed differentiation of
\eqref{eq:ns51-screw} gives, almost everywhere for $t>0$,
\begin{equation}\label{eq:ns51-screw-derivative}
 \mathfrak g'(t)=-4\sinh(t/2)
 +\sum_{n<e^t}\frac{\Lambda(n)}{\sqrt n}
 -\frac12(\psi(1/4)-\log\pi)
 -\frac12e^{-t/2}\Phi(e^{-2t},1,1/4),
\end{equation}
and odd reflection gives its value for $t<0$.
In particular $\mathfrak g'\in L^2_{\rm loc}(\mathbb R)$.
Translation continuity in local $L^2$ and Cauchy--Schwarz show that
$U_f$ is continuous, locally uniformly in $x$. The distributional
identity $F_f'=U_f$ then proves $F_f\in C^1$.

For $f\in V_a$ and $\varphi\in C_c^\infty(I_a)$, the core identity
and symmetry give, in the first-slot-linear convention,
\[
 q_a(f,\varphi)=\langle f,A_a\varphi\rangle
 =-\int_{I_a}F_f(x)\overline{\varphi''(x)}\,dx
 =\int_{I_a}U_f(x)\overline{\varphi'(x)}\,dx.
\]
A nullvector therefore has $F_f''=0$ as a distribution on $I_a$,
which is equivalent to the displayed affine and constant conditions.
Conversely these conditions annihilate $q_a(f,\varphi)$ on the smooth
form core. The form norm, after a fixed-window positive shift, makes
$q_a(f,\cdot)$ continuous. Core density therefore gives
$q_a(f,v)=0$ for every $v\in V_a$. The representation theorem for
closed forms implies $f\in\mathcal D(A_a)$ and $A_af=0$.
The kernels $\mathfrak g(x-y)$ and $\mathfrak g'(x-y)$ are square
integrable on $I_a^2$; for the latter the squared kernel norm is
\[
 \int_{-2a}^{2a}(2a-|t|)|\mathfrak g'(t)|^2\,dt<\infty.
\]
Thus their integral operators are Hilbert--Schmidt, and bounded
projections preserve compactness. Evenness of $\mathfrak g$ gives
the two parity conclusions.
\end{proof}

\paragraph{What changes in the derivative reduction.}
For $f\in H_0^1(I_a)$ one can write the familiar equation
$G_a Df=0$ with $D=i\,d/dx$ and the ordinary mean-zero screw operator
$G_a$. For general $f\in V_a$, $D(E_af)$ is only guaranteed to be a
compactly supported $H^{-1}$ distribution. Integration by parts in
this distributional sense gives
\[
 \mathfrak g*D(E_af)=i\,\mathfrak g'*E_af=iU_f.
\]
This is the extension on this particular derivative range; it does
not assert a bounded extension of $G_a$ to every $H^{-1}$ distribution.
It retains any boundary contributions of zero extension. Therefore
$\mathcal T_af=0$ is meaningful without making $Df$ an ordinary
$L^2$ vector. Suzuki's shifted derivative-completion and generalized
eigenvalue construction already addresses this distinction
\cite[Sections 8.3--8.5]{Suzuki2026}; the proposition gives its
zero-eigenvalue question as a local continuous integral equation.
It is a derived formulation, not a newly discovered RH equivalence.

\begin{proposition}[A zero eigenvalue does not supply generic $H_0^1$ regularity]
\label{prop:ns51-no-generic-bootstrap}
There is an interval $I_r$, a bounded real self-adjoint rank-one
operator $K$ on $L^2(I_r)$, and a nonzero even
$\tau\in\mathcal D(\tfrac12L_\Delta^{\rm Dir})$ such that
\[
 (\tfrac12L_\Delta^{\rm Dir}+K)\tau=0,
 \qquad \tau\notin H_0^1(I_r).
\]
The operator $K$ is bounded on every $L^p(I_r)$, $1\le p\le\infty$,
and $\tau$ has the continuous zero extension and logarithmic boundary
rate appearing in Theorem~\ref{thm:v118-eigen-endpoint-zero}.
This is a countermodel for inference from the principal operator and
bounded perturbation alone, not a counterexample for the arithmetic
$\mathcal K_\lambda$.
\end{proposition}
\begin{proof}
Choose $0<r<e^{-\gamma_{\rm E}}$. The one-dimensional specialization
of \cite[Theorem 1.2]{HLSLogBoundary} supplies the positive even torsion
function $\tau$ with $L_\Delta^{\rm Dir}\tau=1$ and two-sided bounds
$\tau(x)\asymp\ell(r-|x|)^{1/2}$. Its weak equation with right side
$1\in L^2(I_r)$ puts it in the operator domain. Set
\[
 M=\int_{I_r}\tau(x)\,dx>0,\qquad
 Kf=-\frac{1}{2M}\left(\int_{I_r}f(x)\,dx\right)\mathbf1_{I_r}.
\]
This is bounded, real, self-adjoint and rank one, and $K\tau=-1/2$.
If $\tau\in H_0^1(I_r)$, absolute continuity and Cauchy--Schwarz
would give $|\tau(x)|\le\sqrt{r-|x|}\,\|\tau'\|_2$ near either
endpoint. This contradicts its logarithmic lower bound. The two-sided
boundary estimate also supplies the stated continuous extension.
\end{proof}

\paragraph{Named open input: affine-potential injectivity (API).}
For every finite $a>0$, require
\begin{equation}\label{eq:ns51-api}
 f\in V_a,\quad (\mathfrak g*E_af)|_{I_a}\in
 \operatorname{span}\{1,x\}\quad\Longrightarrow\quad f=0.
\end{equation}
The parity versions must use the constants specified above. No
injectivity estimate for the actual arithmetic kernel is proved here.
Injectivity on all of $L^2(I_a)$ would be sufficient but is a stronger
statement than the one established as equivalent to the physical
kernel condition; no regularity upgrade is being assumed.

\begin{proposition}[Sign continuation with an explicit missing input]
\label{prop:ns51-sign-continuation}
Assume the complete fixed-window realizations above, compact resolvent,
continuity of their lowest eigenvalue in $a$, positivity for sufficiently
small $a$, and API for every finite $a>0$. Then every $A_a$ is strictly
positive at fixed $a$. Consequently the complete compact-test Weil
form is nonnegative and Weil's criterion implies RH. No common positive
spectral gap is required. This is a conditional implication; API is open.
\end{proposition}
\begin{proof}
The first four structural inputs, excluding API, are supplied by the
fixed-window theory and \cite[Theorems 1.3--1.4]{Suzuki2026}.
If the continuous lowest eigenvalue ever became nonpositive, it would
be zero at some finite intermediate window. Compact resolvent makes
that lower edge an eigenvalue. Proposition~\ref{prop:ns51-affine-potential}
then contradicts API. Every smooth compact test lies in a finite window,
so its energy is nonnegative. Apply the classical Weil criterion.
\end{proof}

\paragraph{Scope of the finding.}
The unjustified $H_0^1$ shortcut has been replaced by a valid weak
integral formulation. The generic bootstrap is false, while its truth
for a hypothetical arithmetic nullvector has not been decided. The
remaining obstacle is uniqueness for this specific signed arithmetic
integral equation, not a proved negative direction and not a required
uniform inverse bound. Complete even and odd sectors have both been
retained; no source-complement argument is substituted. Growing
near-zero blocks are consistent with fixed-window injectivity.
